% tate-twocolumn.tex — Pandoc カスタムテンプレート（縦書き二段組向け）
%
% 出典: `pandoc -D latex` を基に、以下を調整した派生テンプレート。
%   1. `\maketitle` を `\twocolumn[\@maketitle]` 相当の段抜けタイトルへ差し替え
%      （二段組でタイトルが左段に潰れるのを防ぐ）
%   2. `\documentclass` のデフォルトを ltjtarticle / twocolumn に寄せるが、
%      変数スロットは Pandoc 標準をそのまま維持（defaults file / frontmatter で上書き可能）
%
% 使用方法（ADR-007 defaults 方式）:
%   defaults file の template 項に本ファイルを指定する（path は defaults file 自身の
%   ディレクトリを指す Pandoc の特殊変数を用いて相対指定可能）。
%   クラス・オプションは defaults file の variables で制御:
%     variables:
%       documentclass: ltjtarticle
%       classoption: "twocolumn, a4paper"
%
% 注意: MdTex 固有レイヤ（Obsidian 記法処理・callout/mermaid の Lua フィルタ・
%   codelisting 環境の補完）は、MdTex 側が常時上乗せするため、本テンプレート側では
%   重複定義しない。パッケージの追加が必要な場合は preamble.tex（include-in-header）で。

$passoptions.latex()$
\documentclass[
$for(babel-otherlangs)$
  $babel-otherlangs$,
$endfor$
$if(babel-lang)$
  $babel-lang$,
$endif$
$if(fontsize)$
  $fontsize$,
$endif$
$if(papersize)$
  $papersize$paper,
$endif$
$for(classoption)$
  $classoption$$sep$,
$endfor$
]{$documentclass$}
$if(beamerarticle)$
\usepackage{beamerarticle} % needs to be loaded first
$endif$
\usepackage{xcolor}
$if(geometry)$
\usepackage[$for(geometry)$$geometry$$sep$,$endfor$]{geometry}
$endif$
\usepackage{amsmath,amssymb}
$--
$-- section numbering
$--
$if(numbersections)$
\setcounter{secnumdepth}{$if(secnumdepth)$$secnumdepth$$else$5$endif$}
$else$
\setcounter{secnumdepth}{-\maxdimen} % remove section numbering
$endif$
$fonts.latex()$
$font-settings.latex()$
$common.latex()$
$for(header-includes)$
$header-includes$
$endfor$
$after-header-includes.latex()$
$hypersetup.latex()$

$if(title)$
\title{$title$$if(thanks)$\thanks{$thanks$}$endif$}
$endif$
$if(subtitle)$
\usepackage{etoolbox}
\makeatletter
\providecommand{\subtitle}[1]{% add subtitle to \maketitle
  \apptocmd{\@title}{\par {\large #1 \par}}{}{}
}
\makeatother
\subtitle{$subtitle$}
$endif$
\author{$for(author)$$author$$sep$ \and $endfor$}
\date{$date$}

% --- 段抜けタイトルの準備（二段組向けの調整ポイント） ---
% 二段組クラスオプション（twocolumn）時、\maketitle を段抜きで出す。
% ltjtarticle は \maketitle 内部で \@maketitle を呼ぶ構造なので、
% フロントマター相当の位置で \twocolumn[\@maketitle] を呼び出す。
\makeatletter
\if@twocolumn
  \renewcommand{\maketitle}{%
    \twocolumn[\@maketitle]%
  }
\fi
\makeatother

\begin{document}
$if(has-frontmatter)$
\frontmatter
$endif$
$if(title)$
\maketitle
$if(abstract)$
\begin{abstract}
$abstract$
\end{abstract}
$endif$
$endif$

$for(include-before)$
$include-before$

$endfor$
$if(toc)$
$if(toc-title)$
\renewcommand*\contentsname{$toc-title$}
$endif$
{
$if(colorlinks)$
\hypersetup{linkcolor=$if(toccolor)$$toccolor$$else$$endif$}
$endif$
\setcounter{tocdepth}{$toc-depth$}
\tableofcontents
}
$endif$
$if(lof)$
\listoffigures
$endif$
$if(lot)$
\listoftables
$endif$
$if(linestretch)$
\setstretch{$linestretch$}
$endif$
$if(has-frontmatter)$
\mainmatter
$endif$
$body$

$if(has-frontmatter)$
\backmatter
$endif$
$if(nocite-ids)$
\nocite{$for(nocite-ids)$$it$$sep$, $endfor$}
$endif$
$if(natbib)$
$if(bibliography)$
$if(biblio-title)$
$if(has-chapters)$
\renewcommand\bibname{$biblio-title$}
$else$
\renewcommand\refname{$biblio-title$}
$endif$
$endif$
\bibliography{$for(bibliography)$$bibliography$$sep$,$endfor$}

$endif$
$endif$
$if(biblatex)$
\printbibliography$if(biblio-title)$[title=$biblio-title$]$endif$

$endif$
$for(include-after)$
$include-after$

$endfor$
\end{document}
